../cictikz/src/cictikz/data/tex/cictikz_preamble.tex

% Active cascode (gain boosted) current mirror: amplifiers regulate the
% drains of the mirror devices by driving the cascode gates, multiplying
% the output resistance by the amplifier gain A.
%
% Watch the polarity: the regulated node goes to the amplifier's
% INVERTING input and the reference to the non-inverting one. That way
% the node rising drives the cascode gate down, the cascode conducts
% less, and the node comes back. Wire it the other way and the same
% loop runs away instead of regulating.

\begin{circuitikz}[american, thick, transform shape, circuit ee IEC]

  ../cictikz/src/cictikz/data/tex/cictikz_lib.tex

  \def\xl{0}
  \def\xr{6.4}
  \def\ycas{1.6}    % cascode row bottom
  \def\ytop{3.2}    % top of the cascodes

  % ---- mirror devices ----
  \draw (\xl,0) \vground;
  \draw (\xl,0) \lvnmos{M1}{};
  \draw (\xr,0) \vground;
  \draw (\xr,0) \lvmnmos{M2}{};

  % ---- cascodes, gates facing the amplifiers in the middle ----
  \draw (\xl,\ycas) \lvmnmos{M3}{};
  \draw (\xr,\ycas) \lvnmos{M4}{};

  % ---- input and output branches ----
  \draw (\xl,\ytop) to [short,-o, i_<=$i_{in}$] (\xl,4.4);
  \draw (\xr,\ytop) to [short,-o, i^<=$i_{out}$] (\xr,4.4);

  % ---- mirror gate wire, routed under the ground row ----
  \coordinate (g1) at (M1.G);
  \coordinate (g2) at (M2.G);
  \draw (g1) -- ++(-0.5,0) coordinate (g1a);
  \draw (g1a) -- (g1a |- 0,-0.8) coordinate (g1b);
  \draw (g2) -- ++(0.5,0) coordinate (g2a);
  \draw (g2a) -- (g2a |- 0,-0.8) coordinate (g2b);
  \draw (g1b) -- (g2b);

  % ---- diode connection: M1 gate up to the input node ----
  \fill (g1a) circle (0.075);
  \draw (g1a) -- (g1a |- 0,4.1) coordinate (g1c);
  \draw (g1c) -- (\xl,4.1);
  \fill (\xl,4.1) circle (0.075);

  % ---- left boost amplifier: drives M3's gate ----
  \draw (2.55,2.9) -- (2.55,1.9) -- (1.65,2.4) -- cycle;
  \node[anchor=east, scale=0.8] at (2.5,2.15) {$-$};
  \node[anchor=east, scale=0.8] at (2.5,2.65) {$+$};
  \coordinate (g3) at (M3.G);
  \draw (1.65,2.4) -- (1.15,2.4);
  \draw (1.15,2.4) |- (g3);
  % inverting input senses the drain of M1
  \draw (2.55,2.15) -- (2.95,2.15) -- (2.95,\ycas) -- (\xl,\ycas);
  \fill (\xl,\ycas) circle (0.075);

  % ---- right boost amplifier: drives M4's gate ----
  \draw (3.85,2.9) -- (3.85,1.9) -- (4.75,2.4) -- cycle;
  \node[anchor=west, scale=0.8] at (3.9,2.15) {$-$};
  \node[anchor=west, scale=0.8] at (3.9,2.65) {$+$};
  \coordinate (g4) at (M4.G);
  \draw (4.75,2.4) -- (5.15,2.4);
  \draw (5.15,2.4) |- (g4);
  % inverting input senses the drain of M2
  \draw (3.85,2.15) -- (3.45,2.15) -- (3.45,\ycas) -- (\xr,\ycas);
  \fill (\xr,\ycas) circle (0.075);

  % ---- shared reference for the two non-inverting inputs ----
  \draw (2.55,2.65) -- (3.85,2.65);
  \fill (3.2,2.65) circle (0.075);
  \draw (3.2,2.65) -- (3.2,3.9);
  \draw (3.2,3.9) to [short,-o] ++(0,0.3) node[anchor=south] {$V_{B}$};

  % ---- device names ----
  \node[anchor=west, scale=0.8] at (0.35,0.8) {M1};
  \node[anchor=east, scale=0.8] at (6.05,0.8) {M2};
  \node[anchor=east, scale=0.8] at (-0.35,2.4) {M3};
  \node[anchor=west, scale=0.8] at (6.75,2.4) {M4};

\end{circuitikz}
\end{document}
